% !PRINTED 247
\weberpage{247}
One may call the system $\mathfrak{C}$ a \gesperrt{finite field}, since it contains only finitely many numbers and has the characteristic property of a field: the arithmetic operations can be performed without restriction, except for division by zero (Bd.\,I, \S.\,139)\textsuperscript{1)}. We shall therefore call $\mathfrak{C}$ a \gesperrt{congruence field}.

Let $n$ be called the \gesperrt{degree} and $p$ the \gesperrt{modulus} of the field $\mathfrak{C}$.

If $\alpha$ is a number in $\mathfrak{C}$ different from zero, then the product $\alpha\xi = \eta$ runs through the full number system $\mathfrak{C}$ together with $\xi$. If we exclude $\xi = 0$, then $\eta = 0$ also does not occur, and the product $\Pi$ of all the $\xi$ agrees with the product of all the $\eta$ and is different from zero. But by multiplication of all the equations $\alpha\xi = \eta$ it follows that

\[
\alpha^{p^n-1}\Pi = \Pi,
\]

and after cancelling the common factor $\Pi$ there results

\gesperrt{5. Fermat's theorem:}

\weq{3}{\alpha^{p^n-1} = 1.}

If one multiplies by $\alpha$, one obtains this theorem in the form

\weq{4}{\alpha^{p^n} = \alpha,}

in which it still holds also for $\alpha = 0$. If $\alpha$ is different from zero, then on account of (3) there is a \gesperrt{least} positive exponent $e$ for which

\weq{5}{\alpha^e = 1}

holds, and for which consequently the powers $1, \alpha, \alpha^2, \ldots, \alpha^{e-1}$ are different from one another. One then says that $\alpha$ \gesperrt{belongs to the exponent} $e$. This exponent $e$ must be a divisor of $p^n - 1$. For if $\alpha^m = 1$ for any exponent $m$, we set $m = qe + e'$, where $q$ is an integer and $e' < e$. Then $\alpha^{e'} = 1$ as well, and consequently $e' = 0$ must hold, i.\,e. $m$ must be divisible by $e$. Among these exponents $m$ there is found also $p^n - 1$. If we therefore set

\[
p^n - 1 = ef,
\]

then $\alpha^h$ will always belong to the exponent $e$ whenever $h$ is relatively prime to $e$. For then and only then is $he$ the least positive multiple of $h$ divisible by $e$.

\textsuperscript{1)} See the author's memoir: Die allgemeinen Grundlagen der \gesperrt{Galois'schen} Gleichungstheorie. Mathematische Annalen, Bd.\,43.
